\section{From a growing periodic mode to finite-energy nonuniqueness}
\label{sec:localization}

\ifdefined\CompleteCandidate
We now use the exact background and growing mode from
Sections~\ref{short-bg:section} and~\ref{sec:spectral-assembly}. This is the instability
and localization mechanism of Jia and \v Sver\'ak~\cite{jia-sverak2015},
adapted here to a profile periodic in similarity time. A divergence-free
cutoff localizes the common DSS datum while preserving it near the
origin. A common heat correction and a family of nonlinear correctors then
produce the finite-energy solutions. The periodic right inverse includes
every faster unstable mode; no maximal-growth assumption is added.
\fi

We localize the initial datum and correct the background and its growing
linear perturbation directly. The correction of the common heat tail
decays as $e^{\tau/2}$, while the quadratic mode has size
$e^{2\sigma\tau}$. A periodic spectral splitting at a rate
$\sigma<\beta<2\sigma<1/2$ gives a right inverse below both forcing
rates, without an upper bound on other unstable exponents.

\begin{proposition}\label{short-loc:reduction}
\ifdefined\CompleteCandidate
The background and mode in Proposition~\ref{prop:unstable-background}
imply Theorem~\ref{thm:finite-energy-target}.
\else
Let $\nu,P>0$. Suppose a nonzero real divergence-free field $a$, smooth away
from zero, is axisymmetric without swirl and satisfies
\[
 a(e^{P/2}x)=e^{-P/2}a(x),\qquad
 |D^ja(x)|\le C_j|x|^{-1-j}\quad(0\le j\le4).
\]
Put
\[
 H_\nu[a](y,\tau)=e^{\tau/2}
       [e^{\nu e^\tau\Delta}a](e^{\tau/2}y).
\]
Suppose $U=H_\nu[a]+Q$, $Q\in C_{\rm per}(\mathbb R;H^2(\mathbb R^3))$,
is a real divergence-free $P$-periodic distributional solution of
\[
 \partial_\tau U=\nu\Delta U+\tfrac12y\cdot\nabla U+\tfrac12U
       -\mathbb P\operatorname{div}(U\otimes U),
\]
also axisymmetric without swirl. Suppose, for some $0<\sigma<1/4$,
there is a nonzero real $P$-periodic divergence-free field $b$ of the same
symmetry satisfying
\[
 b,\nabla b\in L^2([0,P]\times\mathbb R^3),\qquad
 (\partial_\tau+\sigma)b
 =\nu\Delta b+\tfrac12y\cdot\nabla b+\tfrac12b
   -\mathbb P\operatorname{div}(U\otimes b+b\otimes U).
\]
Then there are $T>0$, constants $0<c_*\le C_*<\infty$, a nonzero
compactly supported real $u_0\in L^2_{\rm div}(\mathbb R^3)$, and
axisymmetric no-swirl solutions $(u^{(\alpha)},p^{(\alpha)})$,
$\alpha\in[0,1]$, of the unforced viscosity-one Navier--Stokes equation
with this common datum. The datum is smooth away from zero and belongs
to $L^{3,\infty}\cap\bigcap_{1\le q<3}L^q$, but not to $L^3$.
Here $\|f\|_{3,\infty}=\sup_{s>0}s|\{|f|>s\}|^{1/3}$.
Every branch satisfies
\[
 \begin{gathered}
 u^{(\alpha)}\in C([0,T];L^2_{\rm div})\cap L^2(0,T;H^1),\qquad
 u^{(\alpha)}\in C([\delta,T];H^2)\quad(0<\delta<T),\\
 p^{(\alpha)}=\sum_{i,k}R_iR_k(u^{(\alpha)}_iu^{(\alpha)}_k)
       \in L^{5/3}(\mathbb R^3\times(0,T)).
 \end{gathered}
\]
They have the strong $L^2$ trace $u_0$, satisfy the distributional
equation including the initial term $\int u_0\cdot\varphi(0)$,
are smooth at positive times, and obey
\[
 \frac12\|u^{(\alpha)}(t_1)\|_2^2+
 \int_{t_0}^{t_1}\|\nabla u^{(\alpha)}\|_2^2dt
 =\frac12\|u^{(\alpha)}(t_0)\|_2^2
 \quad(0\le t_0\le t_1\le T).
\]
They also satisfy the local energy equality
\eqref{short-loc:local-energy}, including the initial term, and hence
are suitable. The family is uniformly bounded in
$L^\infty(0,T;L^{3,\infty})$ and, for each $1\le r\le\infty$,
$2\le q\le\infty$ with $2/r+3/q>1$, in $L^r(0,T;L^q)$.
For every $\alpha,\alpha'\in[0,1]$ and $0<t\le T$,
\[
 c_*|\alpha-\alpha'|t^{\sigma+1/4}
 \le\|u^{(\alpha)}(t)-u^{(\alpha')}(t)\|_2
 \le C_*|\alpha-\alpha'|t^{\sigma+1/4}.
\]
Consequently their squared $L^2$ distance has time integral at least
$c_*^2|\alpha-\alpha'|^2T^{2\sigma+3/2}/(2\sigma+3/2)$.
\fi
\end{proposition}

\ifdefined\CompleteCandidate
\begin{proof}[Proof of Proposition~\ref{short-loc:reduction} and Theorem~A]
The background construction and the return argument prove
Proposition~\ref{prop:unstable-background}. It remains to localize these
fields and construct the nonlinear solutions. Write $a=u_{\rm dss}$,
$U=U_*$, $Q=r_{\rm heat}$, $b=b_*$ and $P=\mathcal T$.
\else
\begin{proof}
\leavevmode
\fi
\paragraph{Periodic evolution and regularity of the mode.}
The map $D_\nu f(y)=\nu^{-1/2}f(\sqrt \nu\,y)$ changes the profile viscosity
to one: it commutes with the dilation drift and with $\mathbb P$,
$\Delta D_\nu f=D_\nu(\nu\Delta f)$, and
$\operatorname{div}(D_\nu f\otimes D_\nu g)
=D_\nu\operatorname{div}(f\otimes g)$. Also
$H_1[D_\nu a]=D_\nu H_\nu[a]$ and
\[
 \|\nabla^kD_\nu f\|_q=\nu^{(k-1)/2-3/(2q)}\|\nabla^kf\|_q.
\]
We apply this fixed normalization and retain the notation $a,U,Q,b$.
All constants below may depend on these fixed profiles and on the
spatial cutoff, but not on the terminal time $\tau_0\le0$ or on
$\alpha,\alpha'\in[0,1]$.

Let $X=H^1_{\rm div}(\mathbb R^3)$, $Z=H^2_{\rm div}(\mathbb R^3)$,
and write
\[
 \mathcal A=\Delta+\tfrac12y\cdot\nabla+\tfrac12,
 \qquad \mathcal B(\tau)f
 =-\mathbb P\operatorname{div}(U(\tau)\otimes f+f\otimes U(\tau)),
 \qquad \mathcal L=\mathcal A+\mathcal B.
\]
Gaussian convolution and the bounds on $a$ give
$|D^jH_1[a](y,\tau)|\le C_j\langle y\rangle^{-1-j}$ for $j\le2$,
uniformly in phase. Hence $U\in C_{\rm per}C_0$,
$\nabla U\in C_{\rm per}L^3$, and
\[
 \mathcal B\in C_{\rm per}\mathcal L(X,L^2_{\rm div})
                 \cap C_{\rm per}\mathcal L(Z,X).
\]
For the first bound use
$\|U\|_\infty\|\nabla f\|_2+\|\nabla U\|_3\|f\|_6$;
for the second use the $H^2$ product estimate on $Qf$ and the
bounded derivatives of $H_1[a]$. The Leray multiplier is bounded on
all the Sobolev spaces used here.

The free evolution is
\[
 S(s)f(y)=e^{s/2}[e^{(e^s-1)\Delta}f](e^{s/2}y).
\]
It has $X\to X$ norm at most $Ce^{-s/4}$, and $L^2\to X$ and
$X\to Z$ norms at most $Cs^{-1/2}$ for $0<s\le1$. The large-time
gradient estimate follows by applying the heat gradient bound to $f$,
rather than differentiating $f$. Volterra iteration therefore constructs
the periodic evolution $E(\tau,s)$ of $\mathcal L$, bounded on
$L^2,X,Z$ over bounded time strips, with these same short-time smoothing
bounds. The iteration in the smoothing norm uses the integrable
convolution of two kernels $r^{-1/2}$.

These evolutions also represent the given weak mode. On a finite time
interval set $g(x,\tau)=b(e^{-\tau/2}x,\tau)$ and
$\widetilde U(x,\tau)=U(e^{-\tau/2}x,\tau)$. Then
\[
 g_\tau-e^\tau\Delta g+(\sigma-\tfrac12)g
 +e^{\tau/2}\mathbb P\operatorname{div}
       (\widetilde U\otimes g+g\otimes\widetilde U)=0.
\]
The energy hypotheses imply $g\in L^2H^1$, and this equation gives
$g_\tau\in L^2H^{-1}$. Its form is bounded on $H^1\times H^1$;
Young's inequality makes it coercive after adding $C\|g\|_2^2$,
since $e^\tau$ has a positive minimum on the interval. The Hilbert
triple energy identity gives a continuous $L^2$ trace and uniqueness
from every initial time. Thus
$b(\tau)=e^{-\sigma(\tau-s)}E(\tau,s)b(s)$ whenever $b(s)\in H^1$.
Choosing such an earlier phase and using smoothing gives
\begin{equation}\label{short-loc:mode-regularity}
 b\in C_{\rm per}H^2,\qquad
 m_b:=\min_{0\le\tau\le P}\|b(\tau)\|_2>0.
\end{equation}
Indeed a zero phase would, by forward uniqueness and periodicity,
make the entire mode zero.

\paragraph{A right inverse on the ancient interval.}
We next construct an ancient right inverse, including every faster
unstable mode. Multiplication by $U(\tau)\in C_0$ is compact
$X\to L^2$, by compactness on balls and the small uniform spatial tail.
Consequently $\mathcal B(\tau):X\to H^{-1}$ is compact.
For positive $r$, $S(r)\mathcal B(\tau):X\to X$ is compact; its
alternative factorization through $L^2$ has norm $O(r^{-1/2})$.
The Duhamel formula applied to a weakly null sequence, followed by
dominated convergence with this integrable bound, proves
$E(t,s)-S(t-s)$ compact on $X$. For the monodromy
$M=E(P,0)$, periodicity and the free decay therefore give
\[
 r_{\rm ess}(M)\le e^{-P/4}<1.
\]
The spectrum outside this disk consists of isolated eigenvalues with
finite algebraic multiplicity. Choose
$\sigma<\beta<2\sigma<1/2$ so that $|z|=e^{\beta P}$ avoids the
spectrum. The Riesz projections $P_-$ inside that circle and
$P_+=I-P_-$ are real, $P_+X$ is finite dimensional, and for some
$\beta_-<\beta<\beta_+$,
\[
 \|M^nP_-\|\le Ce^{\beta_-nP},\qquad
 \|(M|_{P_+X})^{-n}\|\le Ce^{-\beta_+nP},\quad n\ge0.
\]
These are the spectral-radius estimates for the two restrictions;
no inverse of the full parabolic evolution is taken.

For $0\le\theta\le P$ put
\[
 F_\theta=E(\theta,0)|_{P_+X},\qquad
 G_\theta=(M|_{P_+X})^{-1}P_+E(P,\theta),\qquad
 \Pi_+(\theta)=F_\theta G_\theta,\quad\Pi_-=I-\Pi_+.
\]
Here $G_\theta F_\theta=I$, and the endpoints agree. Extend the
projections periodically. They intertwine the evolution. On the
finite-dimensional exterior fibers its backward extension is
\[
 E_+(kP+\theta,lP+\phi)\Pi_+(\phi)
       =F_\theta(M|_{P_+X})^{k-l}G_\phi.
\]
One-period smoothing shows $P_+X\subset Z$, since
$P_+=M(M|_{P_+X})^{-1}P_+$. Factoring general intervals into
periods and bounded phase strips now gives
\[
 \begin{aligned}
 \|E(\tau,s)\Pi_-(s)\|_{X\to Z}
   &\le C(1+(\tau-s)^{-1/2})e^{\beta_-(\tau-s)}&& (\tau>s),\\
 \|E_+(\tau,s)\Pi_+(s)\|_{X\to Z}
   &\le Ce^{\beta_+(\tau-s)}&& (\tau\le s).
 \end{aligned}
\]
If $P_+X=0$, the second expression is zero. For continuous functions
on $(-\infty,\tau_0]$ use the norms
$\|f\|_{X_\beta}=\sup e^{-\beta\tau}\|f(\tau)\|_X$ and
$\|f\|_{Z_\beta}=\sup e^{-\beta\tau}\|f(\tau)\|_Z$. Then
\begin{equation}\label{short-loc:right-inverse}
 (\mathcal R_\beta f)(\tau)
 =\int_{-\infty}^{\tau}E(\tau,s)\Pi_-(s)f(s)\,ds
  -\int_\tau^{\tau_0}E_+(\tau,s)\Pi_+(s)f(s)\,ds,
 \qquad
 \|\mathcal R_\beta f\|_{Z_\beta}\le C_\beta\|f\|_{X_\beta}.
\end{equation}
Both weighted kernels are integrable. The evolution law verifies
$(\partial_\tau-\mathcal L)\mathcal R_\beta f=f$ and the terminal
condition $\Pi_+(\tau_0)\mathcal R_\beta f(\tau_0)=0$.

\paragraph{Divergence-free truncation of the datum and the common heat correction.}
To obtain a compact datum, define
\[
 A(x)=-\int_0^1r\,x\times a(rx)\,dr,
 \qquad u_0=\operatorname{curl}(\chi A),\qquad d=u_0-a,
\]
where $\chi$ is a smooth compact radial cutoff equal to one near zero.
The derivative bounds give $|D^jA(x)|\le C_j|x|^{-j}$ and
\[
 \operatorname{curl}[r x\times a(rx)]
       =-\partial_r[r^2a(rx)],\qquad \operatorname{curl}A=a.
\]
The latter identity holds across zero: the omitted spherical boundary
term is $O(r^2)$. Thus $u_0$ is compact, divergence-free and in $L^2$.
It is nonzero, since it equals $a$ near zero and DSS rescales any
nonzero value of $a$ into that neighborhood. It is smooth off zero and
$|u_0(x)|\le C|x|^{-1}$; compact support therefore gives
$u_0\in L^{3,\infty}\cap\bigcap_{1\le q<3}L^q$.
To exclude $L^3$, put $\lambda=e^{P/2}>1$ and
$E=\{1<|x|<\lambda\}$. Nonvanishing, continuity and DSS imply
$0<\int_E|a|^3<\infty$. Each of the disjoint annuli $\lambda^{-n}E$
has this same integral, and $u_0=a$ on all sufficiently small such
annuli. Thus $\int|u_0|^3=\infty$.
The field $d$ is smooth,
vanishes near zero, and satisfies inverse-radius derivative bounds.
Set
\[
 H(x,t)=e^{t\Delta}d(x),\qquad
 K(y,\tau)=e^{\tau/2}H(e^{\tau/2}y,e^\tau).
\]
For $\tau\le0$, $\|K(\tau)\|_{W^{2,\infty}}\le Ce^{\tau/2}$.
Write
\[
 \mathcal N(f)=-\mathbb P\operatorname{div}(f\otimes f),\qquad
 \mathcal C(f,g)=\mathbb P\operatorname{div}(f\otimes g+g\otimes f),
 \qquad G=-\mathcal C(U,K)+\mathcal N(K).
\]
Then $\|G(\tau)\|_{H^1}\le Ce^{\tau/2}$. For the $Q$ contribution
this follows from the $H^2$ product bound. For the other contributions,
put $H_a=e^{t\Delta}a$. The physical force
\[
 -\mathbb P\operatorname{div}(H_a\otimes H+H\otimes H_a+H\otimes H)
\]
is uniformly bounded in $H^1$ for $0<t\le1$: on an inner ball the
heat evolution of $d$ and its required derivatives are exponentially
small, whereas derivatives of $H_a$ lose only powers of $t$; outside
that ball the unprojected divergence and its derivative decay as
$\langle x\rangle^{-3}$ and $\langle x\rangle^{-4}$. Apply
$\mathbb P$ after these bounds. Multiplication by $t^{3/2}$ and
evaluation at $\sqrt t\,y$ give $L^2$ and gradient norms bounded by
$Ct^{3/4}$ and $Ct^{5/4}$, which prove the claim for $G$.

\paragraph{Nonlinear correctors and their amplitude dependence.}
Let $v=e^{\sigma\tau}b$. For $\alpha\in[0,1]$ solve in $Z_\beta$
\begin{equation}\label{short-loc:fixed-point}
 z_\alpha=\mathcal R_\beta\{G+\alpha^2\mathcal N(v)
       -\alpha\mathcal C(K,v)
       -\mathcal C(K+\alpha v,z_\alpha)+\mathcal N(z_\alpha)\}.
\end{equation}
The constant, linear and bilinear parts of this map have respective
norms bounded, uniformly in $\alpha$, by
\[
 C_0(e^{(1/2-\beta)\tau_0}+e^{(2\sigma-\beta)\tau_0}),\qquad
 C(e^{\tau_0/2}+e^{\sigma\tau_0}),\qquad Ce^{\beta\tau_0}.
\]
These follow from \eqref{short-loc:right-inverse}, the $H^2$ product
estimate and the bound on $K$. The mixed constant term has the smaller
size $Ce^{(1/2+\sigma-\beta)\tau_0}$. For sufficiently negative
$\tau_0$, all these maps preserve the ball of radius
\[
 \rho=2C_0(e^{(1/2-\beta)\tau_0}+e^{(2\sigma-\beta)\tau_0})
\]
and have Lipschitz constants at most $1/2$. Indeed these constants are
bounded by the displayed linear bound plus $2C\rho e^{\beta\tau_0}$,
which tends to zero. The resulting real $z_\alpha$ satisfy
$\|z_\alpha\|_{Z_\beta}\le\rho$.
Expanding the quadratic terms in \eqref{short-loc:fixed-point} proves
that $U+K+\alpha v+z_\alpha$ solves the exact profile equation.

Subtract the fixed-point equations at $\alpha$ and $\alpha'$.
The terms involving $z_\alpha-z_{\alpha'}$ are absorbed by the
common contraction bound, while $|\alpha^2-\alpha'^2|\le2|\alpha-\alpha'|$.
The remaining terms give
\begin{equation}\label{short-loc:amplitude}
 \|z_\alpha-z_{\alpha'}\|_{Z_\beta}
 \le L_{\tau_0}|\alpha-\alpha'|,\qquad
 L_{\tau_0}=C\bigl(e^{(2\sigma-\beta)\tau_0}
       +e^{(1/2+\sigma-\beta)\tau_0}+\rho e^{\sigma\tau_0}\bigr).
\end{equation}
The three contributions come from $\mathcal N(v)$,
$\mathcal C(K,v)$ and $\mathcal C(v,z_{\alpha'})$, respectively.
Since
\[
 L_{\tau_0}e^{(\beta-\sigma)\tau_0}
 \le C\bigl(e^{\sigma\tau_0}+e^{\tau_0/2}
                    +\rho e^{\beta\tau_0}\bigr)\longrightarrow0,
\]
choose one $\tau_0$ so that this quantity is at most $m_b/2$ and
the preceding contraction bounds hold. This choice is independent
of both amplitudes.

\paragraph{Finite energy, the common trace, and separation.}
With $T=e^{\tau_0}$, the physical solutions are
\begin{equation}\label{short-loc:finite-energy}
 u^{(\alpha)}(x,t)=e^{t\Delta}u_0(x)
       +t^{-1/2}[Q+\alpha v+z_\alpha](x/\sqrt t,\log t).
\end{equation}
The $L^2$ and gradient norms of $t^{-1/2}f(x/\sqrt t)$ are
$t^{1/4}\|f\|_2$ and $t^{-1/4}\|\nabla f\|_2$. Thus the three
remainders have $L^2$ orders $t^{1/4}$, $t^{\sigma+1/4}$ and
$t^{\beta+1/4}$, and squared gradient orders $t^{-1/2}$,
$t^{2\sigma-1/2}$ and $t^{2\beta-1/2}$. All dissipation powers are
integrable. Heat energy and continuity yield the strong trace,
$C_tL^2\cap L_t^2H^1$ and positive-time $C_tH^2$.
The projected equation then gives $\partial_tu^{(\alpha)}\in C_tL^2$
on every positive closed interval. Pairing it with $u^{(\alpha)}$ gives
the energy equality there; strong trace and finite dissipation extend
it to $t_0=0$. For $s\ge2$, the nonlinearity maps $H^s$ to $H^{s-1}$,
and heat smoothing from $H^{s-1}$ to $H^{s+1/2}$ has the integrable
kernel $(t-r)^{-3/4}$. Iteration on smaller positive-time intervals
gives all spatial derivatives; the equation gives all time derivatives.

The interpolation inequality
$\|u\|_{10/3}^{10/3}\le C\|u\|_2^{4/3}\|\nabla u\|_2^2$
gives $u^{(\alpha)}\in L^{10/3}_{x,t}$. The Riesz-transform bounds yield
the stated $L^{5/3}$ pressure and recover the full equation.
Integrating from a positive time and using the strong trace gives
the distributional formulation including $u_0$.

Orthogonal maps fixing the axis commute with the heat and
Leray operators and the nonlinearity. They therefore commute with
the evolution, its spectral projections, and $\mathcal R_\beta$.
The potential $A$ is azimuthal, so $\operatorname{curl}(\chi A)$ is
axisymmetric without swirl. Uniqueness in the common contraction ball
preserves this symmetry for every $z_\alpha$ and
\eqref{short-loc:finite-energy}. By \eqref{short-loc:mode-regularity}
and \eqref{short-loc:amplitude}, for $\tau\le\tau_0$,
\[
 \begin{aligned}
 \| (\alpha-\alpha')v+z_\alpha-z_{\alpha'}\|_2
 &\ge |\alpha-\alpha'|e^{\sigma\tau}
     \bigl(m_b-L_{\tau_0}e^{(\beta-\sigma)\tau_0}\bigr),\\
 \| (\alpha-\alpha')v+z_\alpha-z_{\alpha'}\|_2
 &\le |\alpha-\alpha'|e^{\sigma\tau}
     \bigl(\max_{[0,P]}\|b\|_2+m_b/2\bigr).
 \end{aligned}
\]
Physical scaling gives the two-sided separation with $c_*=m_b/2$;
integrating its square gives the stated integral bound. In particular,
$\alpha\mapsto u^{(\alpha)}-u^{(0)}$ is Lipschitz for the norm
$\sup_{0<t\le T}t^{-\sigma-1/4}\|\cdot(t)\|_2$.

\paragraph{Local energy equality and mixed norms.}
For a fixed branch write $u=u^{(\alpha)}$, $p=p^{(\alpha)}$ and
$e=|u|^2/2$. Its positive-time smoothness gives the classical local
energy identity. For every real
$\phi\in C_c^\infty(\mathbb R^3\times[0,T))$, integration from
$\varepsilon>0$ and then passage to zero give
\ifdefined\CompleteCandidate
\eqref{short-loc:local-energy}.
\else
\begin{equation}\label{short-loc:local-energy}
 \begin{aligned}
 \int_0^T\!\int_{\mathbb R^3}|\nabla u|^2\phi\,dx\,dt
 ={}&\int_0^T\!\int_{\mathbb R^3}
       \bigl[e(\partial_t\phi+\Delta\phi)
                    +(e+p)u\cdot\nabla\phi\bigr]\,dx\,dt\\
 &+\frac12\int_{\mathbb R^3}|u_0|^2\phi(x,0)\,dx.
 \end{aligned}
\end{equation}
\fi
Indeed $|\nabla u|^2$ is integrable; $u\in L^{10/3}_{x,t}$ and
$p\in L^{5/3}_{x,t}$ put both $|u|^3$ and $pu$ in $L^{10/9}_{x,t}$,
hence in $L^1$ on the support of $\phi$. Strong $L^2$
trace implies $|u(\varepsilon)|^2\to|u_0|^2$ in $L^1$.
These facts justify every limit in the identity and, with nonnegative
$\phi$, prove suitability.

Put $F_\alpha=Q+\alpha v+z_\alpha$ and $H_0=e^{t\Delta}u_0$.
The uniform $H^2$ bound on $F_\alpha$ gives, for $2\le q\le\infty$,
\[
 \|t^{-1/2}F_\alpha(x/\sqrt t,\log t)\|_q
 \le C_qt^{(3/q-1)/2}.
\]
The pointwise bound on $u_0$ and Gaussian convolution give
$|H_0(x,t)|\le C/(|x|+\sqrt t)$: rescale by $\sqrt t$ and split
the convolution into $|y-x|<|x|/2$ and its complement when
$|x|>2\sqrt t$. Hence $\|H_0(t)\|_{3,\infty}\le C$ and
$\|H_0(t)\|_\infty\le Ct^{-1/2}$. Since the scaled remainder has
uniform $L^3$ norm, \eqref{short-loc:finite-energy} proves the uniform
weak-$L^3$ assertion.

Heat contraction gives $\|H_0(t)\|_q\le\|u_0\|_q$ for $2\le q<3$.
For $q\ge3$, integrate the distribution function of $|H_0|$,
using its uniform $L^2$ and weak-$L^3$ bounds and its $Ct^{-1/2}$
supremum. Together with the remainder estimate this yields
\[
 \|u^{(\alpha)}(t)\|_q\le C_q
 \begin{cases}
  1,&2\le q<3,\\
  (1+|\log t|)^{1/3},&q=3,\\
  t^{(3/q-1)/2},&3<q\le\infty.
 \end{cases}
\]
For $q=3$ this is integrable to every finite time power. For $q>3$
and finite $r$, the displayed power is integrable exactly when
$r(1-3/q)/2<1$;
for $q<3$ it is uniformly bounded in time. These conditions give
every claimed strict mixed norm, with constants independent of $\alpha$.
\end{proof}
